\documentclass[11pt]{article}
% Required packages
\usepackage[margin=1in]{geometry}
\usepackage{graphicx}
\usepackage{booktabs}
\usepackage{float}
\usepackage{amsmath}
\usepackage{hyperref}
\usepackage{caption}
\usepackage{subcaption}

\title{COS 568: Systems and Machine Learning \\ \Large Project: Learned Index; Task 2, Milestone 1}
\author{}
\date{}

\begin{document}
\maketitle

\section*{Task 2 Milestone 1: HybridPGMLIPP Index}

We introduce \textbf{HybridPGMLIPP}, a hybrid learned index that combines the fast lookup of
LIPP with the efficient insertion of Dynamic PGM (DPGM).
The index is evaluated on the \texttt{fb} dataset (100\,M 64-bit unsigned integers)
on a Della cluster node with Intel CPUs, across two mixed-workload scenarios.

\subsection*{Design}

HybridPGMLIPP maintains two internal structures:
\begin{itemize}
    \item \textbf{LIPP} as the primary read-optimized store, bulk-loaded with the full dataset at build time.
    \item \textbf{DynamicPGM (DPGM)} as a write buffer that absorbs incoming insertions.
\end{itemize}

On each lookup, the index first checks the DPGM write buffer (for recently inserted keys),
then falls through to LIPP. On each insertion, the key is inserted into DPGM.
Once the DPGM buffer accumulates \textbf{100\,000 keys}---equal to $0.1\%$ of the
100\,M-key base dataset---all buffered keys are flushed into LIPP and the buffer is reset.

This flush threshold was chosen to keep the write buffer small relative to the base dataset,
so lookups rarely need to search the buffer, while not flushing so frequently as to incur
excessive LIPP insertion overhead per key.

\paragraph{Flush counts per workload.}
Each workload issues 2\,M operations total:
\begin{itemize}
    \item \textbf{Mixed 10\% insert:} $2\text{M} \times 0.1 = 200{,}000$ insertions $\;\Rightarrow\;$
          $\lfloor 200{,}000 / 100{,}000 \rfloor = \mathbf{2}$ flushes.
    \item \textbf{Mixed 90\% insert:} $2\text{M} \times 0.9 = 1{,}800{,}000$ insertions $\;\Rightarrow\;$
          $\lfloor 1{,}800{,}000 / 100{,}000 \rfloor = \mathbf{18}$ flushes.
\end{itemize}
In both cases the insert count divides evenly into 100\,000, so the buffer is always empty after
the final flush with no residual unflushed keys.

\subsection*{Experimental Results}

Figures~\ref{fig:mix10_tp}--\ref{fig:mix90_sz} show throughput and index size for all three indexes
on the \texttt{fb} dataset under the two mixed workloads.
Table~\ref{tab:results} summarizes the best-variant throughput (mean of 3 runs) for each index.

\begin{table}[H]
\centering
\caption{Best-variant throughput (Mops/s) on \texttt{fb\_100M} mixed workloads.}
\label{tab:results}
\begin{tabular}{lcc}
\toprule
\textbf{Index} & \textbf{Mixed 10\% Insert} & \textbf{Mixed 90\% Insert} \\
\midrule
DynamicPGM     & 1.18 & 4.04 \\
LIPP           & 4.99 & 2.33 \\
HybridPGMLIPP  & 1.41 & 1.51 \\
\bottomrule
\end{tabular}
\end{table}

\begin{figure}[H]
    \centering
    \begin{subfigure}[b]{0.48\textwidth}
        \includegraphics[width=\textwidth]{m2_results/mixed_10insert_throughput.png}
        \caption{Throughput --- Mixed 10\% insert.}
        \label{fig:mix10_tp}
    \end{subfigure}
    \hfill
    \begin{subfigure}[b]{0.48\textwidth}
        \includegraphics[width=\textwidth]{m2_results/mixed_10insert_index_size.png}
        \caption{Index size --- Mixed 10\% insert.}
        \label{fig:mix10_sz}
    \end{subfigure}

    \vspace{0.5em}

    \begin{subfigure}[b]{0.48\textwidth}
        \includegraphics[width=\textwidth]{m2_results/mixed_90insert_throughput.png}
        \caption{Throughput --- Mixed 90\% insert.}
        \label{fig:mix90_tp}
    \end{subfigure}
    \hfill
    \begin{subfigure}[b]{0.48\textwidth}
        \includegraphics[width=\textwidth]{m2_results/mixed_90insert_index_size.png}
        \caption{Index size --- Mixed 90\% insert.}
        \label{fig:mix90_sz}
    \end{subfigure}
    \caption{Throughput and index size for all three indexes on \texttt{fb\_100M} across mixed workloads.
             The index-size plots include a zoomed inset (top right) highlighting the tiny but measurable
             difference between LIPP and HybridPGMLIPP.}
    \label{fig:all}
\end{figure}

\subsection*{Analysis}

\paragraph{Mixed 10\% insert (lookup-heavy).}
With only 200\,000 insertions (2 flushes), HybridPGMLIPP achieves \textbf{1.41\,Mops/s}---below
LIPP (4.99\,Mops/s) and below DynamicPGM (1.18\,Mops/s) as well.
Despite the small number of flushes, each flush requires scanning all 100\,000 DPGM keys and
inserting them one-by-one into LIPP (a slow, conflict-resolving pointer operation).
Between flushes, lookups must also check the DPGM buffer for up to 99\,999 keys before
falling through to LIPP, adding overhead to the dominant lookup path.
LIPP's direct model-guided slot access gives it a decisive advantage in this regime.

\paragraph{Mixed 90\% insert (insert-heavy).}
With 1\,800\,000 insertions (18 flushes), HybridPGMLIPP reaches \textbf{1.51\,Mops/s}---well
below both DynamicPGM (4.04\,Mops/s) and LIPP (2.33\,Mops/s).
The 18 flush events are expensive: each batch of 100\,000 keys must be individually inserted
into LIPP (rather than bulk-loaded), and the DPGM buffer must also be rebuilt from scratch
after each flush.
DPGM's amortized logarithmic-merge strategy handles high insert volumes far more efficiently.

\paragraph{Index size.}
As shown in Figures~\ref{fig:mix10_sz} and~\ref{fig:mix90_sz}, all three indexes are measured
after completing the full 2\,M-operation workload (including all insertions and flushes).
HybridPGMLIPP's index size is dominated by its LIPP component, which holds all 100\,M keys
after every flush has been committed:
\begin{itemize}
    \item \textbf{Mixed 10\% insert:}
          LIPP\,=\,12.700947\,GB vs.\ HybridPGMLIPP\,=\,12.702649\,GB
          ($\Delta \approx \mathbf{1.70\,MB}$, i.e.\ $0.013\%$).
    \item \textbf{Mixed 90\% insert:}
          LIPP\,=\,12.656663\,GB vs.\ HybridPGMLIPP\,=\,12.658368\,GB
          ($\Delta \approx \mathbf{1.71\,MB}$, i.e.\ $0.014\%$).
\end{itemize}
This tiny overhead arises from the DPGM write buffer: it holds at most 100\,000 keys
at any point in time before a flush.
Each key--value record in DynamicPGM occupies $\sim$16--17 bytes, so the buffer contributes
at most $100{,}000 \times 17 \approx \mathbf{1.7\,MB}$, which matches the measured difference
exactly.
Since the 0.1\% flush threshold is specifically designed to keep the buffer tiny,
this negligible overhead is an inherent design property: a larger threshold would produce a
larger buffer (and more visible overhead) but would hurt lookup performance by forcing
lookups to scan a bigger write buffer.
Because this $\sim$1.7\,MB difference is invisible on a full-scale bar chart, the index-size
plots (Figures~\ref{fig:mix10_sz} and~\ref{fig:mix90_sz}) include a zoomed inset that
magnifies the LIPP vs.\ HybridPGMLIPP bars so the difference can be seen.
DynamicPGM ($\approx$1.71\,GB) remains $\approx$7.4$\times$ smaller than either
LIPP-based structure.

\paragraph{Summary.}
In its current form, HybridPGMLIPP does not outperform any baseline on either workload.
The bottleneck is the flush mechanism: inserting keys into LIPP one-by-one is expensive,
and the DPGM buffer lookup overhead hurts the read path.
Future improvements could include bulk-loading flushed keys into a new LIPP instance via
merge, tuning the flush threshold (larger threshold $\Rightarrow$ fewer but costlier flushes;
smaller $\Rightarrow$ more frequent but cheaper flushes), or adopting a two-level
bulk-load strategy to reduce per-flush LIPP insertion cost.

\end{document}
